%% kapa: exam=1 round=33 year=2022 period=2 form=A subject=appraisal_law subject_ko=감정평가관계법규 seq=33 number=33
%% source: Q-Net 감정평가사 기출문제 (cst00309) artlSeq=5209997 file=raw/1008/5209997/2022년도 제33회 감정평가사 1차시험 2교시(감정평가 관계법규, 회계학).pdf
%% conversion: kordoc 4.24.51 -> tools/gen1.py (auto); review=proofread-sonnet; flags=none
\begin{question}{33}
부동산등기법령상 등기신청 및 등기의 효력발생시기에 관한 설명으로 옳은 것을 모두 고른 것은?
\begin{framedtext}
ㄱ. 소유권보존등기 또는 소유권보존등기의 말소등기는 등기명의인으로 될 자 또는 등기명의인이 단독으로 신청한다. \par
ㄴ. 대표자가 있는 법인 아닌 사단에 속하는 부동산의 등기 신청에 관하여는 그 사단의 대표자를 등기권리자 또는 등기의무자로 한다. \par
ㄷ. 등기신청은 해당 부동산이 다른 부동산과 구별될 수 있게 하는 정보가 전산정보처리조직에 저장된 때 접수된 것으로 본다. \par
ㄹ. 등기관이 등기를 마친 경우 그 등기는 등기완료의 통지를 한 때부터 효력을 발생한다. \par
\end{framedtext}
\begin{choices}
  \item ㄱ, ㄴ
  \item ㄱ, ㄷ
  \item ㄷ, ㄹ
  \item ㄴ, ㄷ, ㄹ
  \item ㄱ, ㄴ, ㄷ, ㄹ
\end{choices}
\end{question}
\answer{2}
